\documentclass[11pt]{article}

% Answers and hints for the practice problems of Lecture 3
\usepackage[margin=1in]{geometry}
\usepackage[T1]{fontenc}
\usepackage{lmodern}
\usepackage{amsmath,amssymb,amsthm}
\usepackage{enumitem}
\usepackage{xr-hyper}
\usepackage{hyperref}
% Problem numbers are read from practice3.aux, so compile practice3.tex first.
\externaldocument{practice3}[practice3.pdf]
\usepackage{xcolor}

\title{Lecture 3: Practice Problems --- Answers and Hints}
\author{Math for Machine Learning}
\date{\today}

\newcommand{\norm}[1]{\lVert #1 \rVert}
\newcommand{\ip}[2]{\langle #1, #2 \rangle}

\DeclareMathOperator{\tr}{tr}
\DeclareMathOperator{\diag}{diag}
\DeclareMathOperator{\Div}{div}
\DeclareMathOperator{\curl}{curl}
\DeclareMathOperator{\softmax}{softmax}
\DeclareMathOperator{\ReLU}{ReLU}

\begin{document}

\maketitle

\section*{Multivariable Functions}

\begin{description}[leftmargin=0pt, itemsep=0.3em]
    \item[\ref{pr:level-sets}.] (a) $\mathbb{R}^2$; the parallel lines $2x - y = c - 1$, with normal vector $(2, -1)$. (b) $\mathbb{R}^2$; for $c > 0$ the ellipses $\frac{x^2}{c/4} + \frac{y^2}{c} = 1$ with half-axes $\frac{\sqrt{c}}{2}$ and $\sqrt{c}$, the origin for $c = 0$, and the empty set for $c < 0$. (c) $\mathbb{R}^2$; the parabolas $y = x^2 + c$, which are translates of each other. (d) $\mathbb{R}^2 \setminus \{(0, 0)\}$; the circle of radius $e^{c/2}$ for every $c \in \mathbb{R}$.
    \item[\ref{pr:limits}.] (a) $0$, since $0 \leq \frac{x^2 y^2}{x^2 + y^2} \leq y^2$. (b) Does not exist: the function is $1$ on the $x$-axis and $-1$ on the $y$-axis. (c) $0$: $|x|^3 \leq |x|(x^2 + y^2)$ and similarly for $y$, so the absolute value is at most $|x| + |y|$. (d) $0$: $|xy| \leq \frac{1}{2}(x^2 + y^2)$, so the absolute value is at most $\frac{1}{2}\sqrt{x^2 + y^2}$.
    \item[\ref{pr:lines-not-enough}.] For $\mathbf{v} = (a, b)$ with $a \neq 0$, $f(ta, tb) = \frac{tab^2}{a^2 + t^2 b^4} \to 0$; for $a = 0$, $f = 0$. On the parabola $x = y^2$, $f(y^2, y) = \frac{y^4}{2y^4} = \frac{1}{2}$ for $y \neq 0$.
    \item[\ref{pr:mse-quadratic}.] (a) With $\sum x_i = 6$, $\sum x_i^2 = 14$, $\sum y_i = 10$, $\sum x_i y_i = 23$, $\sum y_i^2 = 38$: $L = \frac{1}{3}(14w^2 + 12wb + 3b^2 - 46w - 20b + 38)$. (b) $L(\mathbf{p}) = \mathbf{p}^\top A\mathbf{p} - 2\mathbf{c}^\top\mathbf{p} + d$ with $A = \frac{1}{3}\begin{pmatrix} 14 & 6 \\ 6 & 3 \end{pmatrix}$, which is positive definite ($14 > 0$, $\det = \frac{6}{9} > 0$). With $\mathbf{p}^* = A^{-1}\mathbf{c}$, $L(\mathbf{p}) = L(\mathbf{p}^*) + (\mathbf{p} - \mathbf{p}^*)^\top A(\mathbf{p} - \mathbf{p}^*)$, whose level sets are ellipses centered at $\mathbf{p}^*$, with axes along the eigenvectors of $A$.
\end{description}

\section*{Derivatives of Functions of One Variable}

\begin{description}[leftmargin=0pt, itemsep=0.3em]
    \item[\ref{pr:activation-derivs}.] (a) Product rule and $\sigma' = \sigma(1 - \sigma)$. (b) $(\log\sigma)' = \frac{\sigma'}{\sigma} = 1 - \sigma$, and $1 - \sigma(x) = \frac{e^{-x}}{1 + e^{-x}} = \frac{1}{e^x + 1} = \sigma(-x)$. (c) Both one-sided difference quotients tend to $1$ (the left one is $\frac{e^h - 1}{h} \to 1$), so $f'(0) = 1$; $f'(x) = 1$ for $x \geq 0$ and $f'(x) = e^x$ for $x < 0$. The derivative $f'$ is continuous, but not differentiable at $0$: its one-sided derivatives are $e^0 = 1$ (left) and $0$ (right). (d) The right quotient is $1$, the left one $0.1$; not differentiable at $0$.
    \item[\ref{pr:taylor-sqrt}.] (a) $T_3(x) = 1 + \frac{x}{2} - \frac{x^2}{8} + \frac{x^3}{16}$. (b) $T_3(0.1) = 1.0488125$, while $\sqrt{1.1} = 1.0488088\ldots$. (c) $f^{(4)}(x) = -\frac{15}{16}(1 + x)^{-7/2}$, which has absolute value at most $\frac{15}{16}$ for $x \geq 0$, so the error is at most $\frac{15/16}{4!}(0.1)^4 \approx 3.9 \cdot 10^{-6}$ (the actual error is $3.65 \cdot 10^{-6}$).
    \item[\ref{pr:taylor-limits}.] (a) $\sin x - x = -\frac{x^3}{6} + (\text{terms of order } x^5)$, so the limit is $-\frac{1}{6}$. (b) $\frac{1}{2}$. (c) $\frac{1}{2}$. (d) $\log(1 + x) - x = -\frac{x^2}{2} + \frac{x^3}{3} - \cdots$, so the limit is $-\frac{1}{2}$.
    \item[\ref{pr:mvt}.] $f(3) - f(0) = 3f'(\xi) \geq -3$ for some $\xi \in (0, 3)$, so $f(0) \leq f(3) + 3 = 8$. The function $f(x) = 8 - x$ attains this.
    \item[\ref{pr:finite-diff}.] (a) Forward: $1.051709$, error $0.0517$, bound $\frac{e^{0.1}}{2}(0.1) \approx 0.0553$. Central: $1.001668$, error $0.00167$, bound $\frac{e^{0.1}}{6}(0.1)^2 \approx 0.00184$. (b) The errors are $0.0254$ and $0.000417$: they shrink by the factors $2.03$ and $4.00$, since the errors are approximately $\frac{f''(0)}{2}h$ and $\frac{f'''(0)}{6}h^2$.
\end{description}

\section*{Partial Derivatives and Gradients}

\begin{description}[leftmargin=0pt, itemsep=0.3em]
    \item[\ref{pr:gradients-compute}.] (a) $(2xy^3 - 4y, \, 3x^2y^2 - 4x)^\top = (2, -1)^\top$. (b) $\frac{2\mathbf{x}}{\norm{\mathbf{x}}^2} = \frac{1}{9}(2, 4, 4)^\top$. (c) $2(\mathbf{a}^\top\mathbf{x})\mathbf{a}$. (d) $-e^{-\norm{\mathbf{x}}^2/2}\mathbf{x}$. (e) $2A\mathbf{x} + \mathbf{b}$.
    \item[\ref{pr:directional}.] (a) $(6xy, \, 3x^2 - 3y^2)^\top = (12, -9)^\top$. (b) The direction is $\frac{1}{5}(3, 4)^\top$, and the directional derivative is $\frac{1}{5}(36 - 36) = 0$. (c) Fastest increase in the direction $\frac{1}{5}(4, -3)^\top$ with rate $\norm{\nabla f(P)} = 15$; fastest decrease in the direction $\frac{1}{5}(-4, 3)^\top$ with rate $-15$. (d) $\pm\frac{1}{5}(3, 4)^\top$, the directions tangent to the level curve through $P$.
    \item[\ref{pr:temperature}.] (a) $\nabla T(2, 1) = (-2x, -4y)^\top = (-4, -4)^\top$: move in the direction $-\frac{1}{\sqrt{2}}(1, 1)^\top$, towards the origin; the temperature increases at the rate $4\sqrt{2} \approx 5.66$ per unit length. (b) The tangent vectors of $y = \frac{x^2}{4}$ are multiples of $(1, \frac{x}{2})$, and $\nabla T(x, \frac{x^2}{4}) = (-2x, -x^2)^\top = -2x(1, \frac{x}{2})^\top$. The bug ends up at the hottest point, the origin, where $T = 100$.
    \item[\ref{pr:directional-not-diff}.] (a) $\frac{x^2}{x^2 + y^2} \leq 1$, so $|f(x, y)| \leq |y| \to 0$. (b) $\frac{f(ta, tb)}{t} = \frac{a^2 b}{a^2 + b^2} = a^2 b$. (c) Both partial derivatives are $0$ (take $\mathbf{v} = \mathbf{e}_1, \mathbf{e}_2$). If $f$ were differentiable, all directional derivatives would be $\nabla f(0, 0)^\top\mathbf{v} = 0$, but $D_{\mathbf{v}} f(0, 0) = \frac{1}{2\sqrt{2}}$ for $\mathbf{v} = \frac{1}{\sqrt{2}}(1, 1)^\top$.
    \item[\ref{pr:gradient-ip}.] (a) $Df(1, 1)\mathbf{h} = 2h_1 + 2h_2$ and $\ip{\mathbf{g}}{\mathbf{h}}_A = g_1 h_1 + 4g_2 h_2$, so $\mathbf{g} = (2, \frac{1}{2})^\top = A^{-1}\nabla f(1, 1)$. (b) By the Cauchy--Schwarz inequality for $\ip{\cdot}{\cdot}_A$, $\ip{\mathbf{g}}{\mathbf{h}}_A \leq \norm{\mathbf{g}}_A\norm{\mathbf{h}}_A = \sqrt{5}$, with equality only for $\mathbf{h} = \mathbf{g}/\norm{\mathbf{g}}_A$. (c) $\nabla f(1, 1) = (2, 2)^\top$ is the steepest direction for the Euclidean norm. In the norm $\norm{\cdot}_A$, steps in the $x_2$-direction are twice as long, so the steepest direction leans towards $x_1$.
\end{description}

\section*{Jacobians: Derivatives of Vector-Valued Functions}

\begin{description}[leftmargin=0pt, itemsep=0.3em]
    \item[\ref{pr:jacobian-squaring}.] (a) $DF = \begin{pmatrix} 2x & -2y \\ 2y & 2x \end{pmatrix}$, $\det DF = 4(x^2 + y^2)$. (b) $F(1, 2) = (-3, 4)^\top$ and $DF(1, 2) = \begin{pmatrix} 2 & -4 \\ 4 & 2 \end{pmatrix}$, so $F(1.1, 1.9) \approx (-3, 4)^\top + (0.6, 0.2)^\top = (-2.4, 4.2)^\top$; exact: $(-2.4, 4.18)^\top$. (c) By the factor $4(x^2 + y^2) = 4|z|^2$; areas are preserved locally on the circle $x^2 + y^2 = \frac{1}{4}$. (The Jacobian is $2|z|$ times a rotation, so $F$ maps small squares to squares.)
    \item[\ref{pr:jacobians-ml}.] (a) The diagonal matrix with entries $1$ where $x_i > 0$ and $0$ where $x_i < 0$. (b) $Df(\mathbf{x}) = \frac{1}{\norm{\mathbf{x}}}I - \frac{\mathbf{x}\mathbf{x}^\top}{\norm{\mathbf{x}}^3} = \frac{1}{\norm{\mathbf{x}}}(I - \mathbf{u}\mathbf{u}^\top)$, and $(I - \mathbf{u}\mathbf{u}^\top)\mathbf{x} = \mathbf{x} - \mathbf{u}(\mathbf{u}^\top\mathbf{x}) = \mathbf{0}$. Expected, since $f(t\mathbf{x}) = f(\mathbf{x})$ for $t > 0$: $f$ is constant along the ray through $\mathbf{x}$. (c) $\frac{1}{T}\big(\diag(\mathbf{s}) - \mathbf{s}\mathbf{s}^\top\big)$ with $\mathbf{s} = \softmax(\mathbf{z}/T)$, by the chain rule.
    \item[\ref{pr:polar-integrals}.] (a) $\int_0^{2\pi}\int_0^1 r^2 \cdot r \, dr \, d\varphi = \frac{\pi}{2}$. (b) $2\pi\big[-e^{-r^2/2}\big]_0^\infty = 2\pi$. (c) $\frac{1}{2\pi}\int_0^{2\pi}\int_0^r e^{-s^2/2}s \, ds \, d\varphi = 1 - e^{-r^2/2}$, which is $0.95$ for $r = \sqrt{2\log 20} \approx 2.45$. In two dimensions, more of the probability lies far from the center than in one dimension.
\end{description}

\section*{The Chain Rule}

\begin{description}[leftmargin=0pt, itemsep=0.3em]
    \item[\ref{pr:chain-compose}.] $Df(u, v) = (2u, -2)$ and $Dg(x, y) = \begin{pmatrix} 1 & 1 \\ y & x \end{pmatrix}$, so $D(f \circ g) = (2(x + y) - 2y, \, 2(x + y) - 2x) = (2x, 2y)$. Directly: $f(g(x, y)) = (x + y)^2 - 2xy = x^2 + y^2$.
    \item[\ref{pr:loss-gradients}.] (a) With $z = y\,\mathbf{w}^\top\mathbf{x}$: $\frac{d\ell}{dz} = \frac{-e^{-z}}{1 + e^{-z}} = -\sigma(-z)$ and $Dz(\mathbf{w}) = y\,\mathbf{x}^\top$ ($1 \times d$), so $\nabla\ell(\mathbf{w}) = -y\,\sigma(-y\,\mathbf{w}^\top\mathbf{x})\,\mathbf{x}$. (b) $X^\top(X\mathbf{w} - \mathbf{y}) + \lambda\mathbf{w}$, from the row vector $(X\mathbf{w} - \mathbf{y})^\top X$ of size $(1 \times N)(N \times d)$. (c) $\frac{A^\top(A\mathbf{x} - \mathbf{b})}{\norm{A\mathbf{x} - \mathbf{b}}}$, from $\nabla\norm{\mathbf{r}} = \frac{\mathbf{r}}{\norm{\mathbf{r}}}$ and $D(A\mathbf{x} - \mathbf{b}) = A$.
    \item[\ref{pr:curve-temperature}.] (a) $\nabla T = (2xy, x^2)^\top$ and $\mathbf{x}'(t) = (-\sin t, \cos t)^\top$, so $\frac{d}{dt}T = -2\cos t\sin^2 t + \cos^3 t$, which is $\cos t(\cos^2 t - 2\sin^2 t)$ and the derivative of $\cos^2 t \sin t$. At $t = \frac{\pi}{4}$: $-\frac{\sqrt{2}}{4} \approx -0.354$. (b) $\cos t = 0$ (where $T = 0$), or $\tan^2 t = \frac{1}{2}$. The maximum $\frac{2}{3\sqrt{3}} \approx 0.385$ is attained where $\cos^2 t = \frac{2}{3}$ and $\sin t = \frac{1}{\sqrt{3}}$, i.e., at $(\pm\sqrt{2/3}, \frac{1}{\sqrt{3}})$.
    \item[\ref{pr:tangent}.] (a) $\nabla = (2x, 8y) = (4, 8)$: $4(x - 2) + 8(y - 1) = 0$, i.e., $x + 2y = 4$. (b) $x + 2y + 2z = 9$. (c) $\nabla F = (2, 4, -1)$: $2(x - 1) + 4(y - 2) - (z - 5) = 0$, i.e., $z = 2x + 4y - 5$.
    \item[\ref{pr:implicit}.] $2x + y + xy' + 3y^2 y' = 0$; at $(1, 1)$: $3 + 4y' = 0$, so $y'(1) = -\frac{3}{4}$. In general, $\frac{d}{dx}F(x, y(x)) = \frac{\partial F}{\partial x} + \frac{\partial F}{\partial y}y'(x) = 0$ by the chain rule.
\end{description}

\section*{Gradients with Respect to Matrices}

\begin{description}[leftmargin=0pt, itemsep=0.3em]
    \item[\ref{pr:matrix-grads}.] (a) $\tr(A(X + H)^\top) = \tr(AX^\top) + \tr(AH^\top)$, and $\tr(AH^\top) = \sum_{i, j} A_{ij}H_{ij} = \ip{A}{H}$. (b) $\norm{(X + H)\mathbf{a}}^2 = \norm{X\mathbf{a}}^2 + 2(X\mathbf{a})^\top H\mathbf{a} + \norm{H\mathbf{a}}^2$, and $(X\mathbf{a})^\top H\mathbf{a} = \ip{X\mathbf{a}\mathbf{a}^\top}{H}$. (c) $\tr((X + H)^2) = \tr(X^2) + 2\tr(XH) + \tr(H^2)$, and $\tr(XH) = \ip{X^\top}{H}$. (d) Chain rule: $\norm{X}_F = (\norm{X}_F^2)^{1/2}$, so the gradient is $\frac{1}{2\norm{X}_F} \cdot 2X$.
    \item[\ref{pr:regularized-layer}.] (a) The linear layer example of the lecture with $\boldsymbol{\delta} = W\mathbf{x} - \mathbf{y}$, plus $\nabla\frac{\lambda}{2}\norm{W}_F^2 = \lambda W$. (b) $W^*\mathbf{x} = \frac{\norm{\mathbf{x}}^2}{\norm{\mathbf{x}}^2 + \lambda}\mathbf{y}$, so $(W^*\mathbf{x} - \mathbf{y})\mathbf{x}^\top + \lambda W^* = \big(\frac{\norm{\mathbf{x}}^2}{\norm{\mathbf{x}}^2 + \lambda} - 1 + \frac{\lambda}{\norm{\mathbf{x}}^2 + \lambda}\big)\mathbf{y}\mathbf{x}^\top = 0$. (c) $W^*\mathbf{x} \to \mathbf{y}$ as $\lambda \to 0$ (a perfect fit), and $W^* \to 0$ as $\lambda \to \infty$: the regularization shrinks the weights.
    \item[\ref{pr:det-2x2}.] (a) $\nabla\det(X) = \begin{pmatrix} d & -c \\ -b & a \end{pmatrix}$, and $X^{-1} = \frac{1}{\det X}\begin{pmatrix} d & -b \\ -c & a \end{pmatrix}$, so $\det(X)X^{-\top} = \begin{pmatrix} d & -c \\ -b & a \end{pmatrix}$. (b) $\det(X) > 0$, so $\log\det X = \log|\det X|$ has the gradient $X^{-\top} = X^{-1}$, since $X^{-1}$ is symmetric.
    \item[\ref{pr:softmax-temperature}.] (a) $\mathbf{p} \approx (0.090, 0.245, 0.665)^\top$, $\ell \approx 0.408$, $\nabla\ell \approx (0.090, 0.245, -0.335)^\top$. (b) Chain rule with $D(\mathbf{z}/T) = \frac{1}{T}I$. For $T = \frac{1}{2}$: $\mathbf{p} \approx (0.016, 0.117, 0.867)^\top$, $\ell \approx 0.143$, $\nabla\ell \approx (0.032, 0.235, -0.266)^\top$. For $T = 2$: $\mathbf{p} \approx (0.186, 0.307, 0.506)^\top$, $\ell \approx 0.680$, $\nabla\ell \approx (0.093, 0.154, -0.247)^\top$.
\end{description}

\section*{Backpropagation and Automatic Differentiation}

\begin{description}[leftmargin=0pt, itemsep=0.3em]
    \item[\ref{pr:autodiff-graph}.] Values: $v_1 = 2$, $v_2 = 0$, $v_3 = 1$, $f = v_4 = 2$. Forward mode with $(\dot x, \dot y) = (1, 0)$: $\dot v_1 = 1$, $\dot v_2 = \dot x y + x\dot y = 0$, $\dot v_3 = e^{v_2}\dot v_2 = 0$, $\dot v_4 = \dot v_1 v_3 + v_1\dot v_3 = 1$; with $(0, 1)$: $\dot v_1 = 1$, $\dot v_2 = 2$, $\dot v_3 = 2$, $\dot v_4 = 1 + 4 = 5$. Reverse mode: $\bar v_4 = 1$, $\bar v_1 = v_3 = 1$, $\bar v_3 = v_1 = 2$, $\bar v_2 = \bar v_3 e^{v_2} = 2$, $\bar x = \bar v_1 + \bar v_2 y = 1$, $\bar y = \bar v_1 + \bar v_2 x = 5$. So $\nabla f(2, 0) = (1, 5)^\top$; directly, $\frac{\partial f}{\partial x} = e^{xy}(1 + (x + y)y)$ and $\frac{\partial f}{\partial y} = e^{xy}(1 + (x + y)x)$.
    \item[\ref{pr:backprop-numbers}.] (a) $\mathbf{z} = (2, -0.5)^\top$, $\mathbf{a} = (2, 0)^\top$, $\hat y = 2.5$, $L = 0.5$. (b) $\hat y - y = 1$; $\nabla_{\mathbf{w}_2} L = \mathbf{a} = (2, 0)^\top$; $\frac{\partial L}{\partial b_2} = 1$; $\nabla_{\mathbf{a}} L = \mathbf{w}_2 = (1, -1)^\top$; $\nabla_{\mathbf{z}} L = (1, 0)^\top \odot (1, -1)^\top = (1, 0)^\top$; $\nabla_{W_1} L = \begin{pmatrix} 1 & 0.5 \\ 0 & 0 \end{pmatrix}$; $\nabla_{\mathbf{b}_1} L = (1, 0)^\top$. (c) The second hidden unit is inactive ($z_2 < 0$): its incoming weights (the second row of $W_1$ and the second entry of $\mathbf{b}_1$) get zero gradients since $\ReLU'(z_2) = 0$, and its outgoing weight $(\mathbf{w}_2)_2$ gets a zero gradient since $a_2 = 0$. With only this example, these weights never change: a ``dead'' ReLU unit.
    \item[\ref{pr:modes-cost}.] (a) $f: \mathbb{R}^{1000} \to \mathbb{R}$: $1000$ forward passes or $1$ reverse pass. $f: \mathbb{R} \to \mathbb{R}^{1000}$: $1$ forward pass or $1000$ reverse passes. (b) $J_2 J_1$ costs $100 \cdot 100 \cdot 10^4 = 10^8$, and then $J_3(J_2 J_1)$ another $10^6$. $J_3 J_2$ costs $10^4$, and then $(J_3 J_2)J_1$ another $10^6$. So the second order, which is reverse mode, is about $100$ times cheaper.
    \item[\ref{pr:vanishing}.] (a) By the chain rule, $f_L'(x)$ is a product of $L$ values of $\sigma'$, each at most $\frac{1}{4}$. For $L = 10$: $4^{-10} \approx 9.5 \cdot 10^{-7}$. (b) $g_{l+1}'(x) = \sigma'(w\,g_l(x))\,w\,g_l'(x)$, so each layer contributes a factor of absolute value at most $\frac{|w|}{4}$. Gradients can only avoid vanishing if $|w| > 4$, but large weights push $\sigma$ into its flat regions, where $\sigma'$ is tiny; this is one reason why deep networks use ReLU.
\end{description}

\section*{Higher Derivatives and Taylor's Theorem}

\begin{description}[leftmargin=0pt, itemsep=0.3em]
    \item[\ref{pr:hessian-compute}.] (a) $\nabla f = (3x^2 - 2y^2, \, -4xy + 4y^3)^\top$ and $\nabla^2 f = \begin{pmatrix} 6x & -4y \\ -4y & -4x + 12y^2 \end{pmatrix}$; both mixed partial derivatives are $-4y$. (b) $f(1, 1) = 0$, $\nabla f(1, 1) = (1, 0)^\top$, $\nabla^2 f(1, 1) = \begin{pmatrix} 6 & -4 \\ -4 & 8 \end{pmatrix}$, so $f(1 + h_1, 1 + h_2) \approx h_1 + 3h_1^2 - 4h_1h_2 + 4h_2^2$. At $(0.1, -0.1)$: $0.21$; exact: $0.2051$. (c) $\frac{1}{2}(6 - 8 + 8) = 3$. The largest and smallest curvatures are the eigenvalues $7 \pm \sqrt{17}$, i.e., $11.12$ and $2.88$.
    \item[\ref{pr:hessians-ml}.] (a) $X^\top X + \lambda I$. (b) $\frac{1}{\norm{\mathbf{x}}}(I - \mathbf{u}\mathbf{u}^\top)$ with $\mathbf{u} = \frac{\mathbf{x}}{\norm{\mathbf{x}}}$ (the Jacobian of the gradient $\mathbf{u}$, Problem~\ref{pr:jacobians-ml}). Eigenvalue $0$ with eigenvector $\mathbf{x}$, and eigenvalue $\frac{1}{\norm{\mathbf{x}}}$ on the orthogonal complement of $\mathbf{x}$. Along the ray through $\mathbf{x}$, $\norm{t\mathbf{x}} = t\norm{\mathbf{x}}$ is linear, so there is no curvature in this direction. (c) $\sigma(\mathbf{w}^\top\mathbf{x})\big(1 - \sigma(\mathbf{w}^\top\mathbf{x})\big)\mathbf{x}\mathbf{x}^\top$, of rank $1$ (for $\mathbf{x} \neq \mathbf{0}$).
    \item[\ref{pr:lse-taylor}.] (a) $\nabla f(0, 0) = (\frac{1}{2}, \frac{1}{2})^\top$ and $\nabla^2 f(0, 0) = \frac{1}{4}\begin{pmatrix} 1 & -1 \\ -1 & 1 \end{pmatrix}$. (b) $\log 2 + \frac{1}{2}(0.2 - 0.1) + \frac{1}{8}(0.2 + 0.1)^2 \approx 0.754397$; exact: $0.754355$. (c) Eigenvalue $\frac{1}{2}$ with eigenvector $(1, -1)$, and $0$ with eigenvector $(1, 1)$: along $(1, 1)$, $f(t, t) = \log 2 + t$ is linear.
    \item[\ref{pr:schwarz-counterexample}.] (a) $f(x, y) = x\,g(x, y)$ with $g = \frac{y(x^2 - y^2)}{x^2 + y^2}$, and $f(0, y) = 0$, so $\frac{\partial f}{\partial x}(0, y) = \lim_{h \to 0}g(h, y) = -y$ (also for $y = 0$). Similarly $f = y\,k$ with $k = \frac{x(x^2 - y^2)}{x^2 + y^2}$ gives $\frac{\partial f}{\partial y}(x, 0) = x$. (b) Differentiate $-y$ with respect to $y$, and $x$ with respect to $x$. Schwarz's theorem needs continuous second partial derivatives, and they are not continuous at the origin.
\end{description}

\section*{Unconstrained Optimization}

\begin{description}[leftmargin=0pt, itemsep=0.3em]
    \item[\ref{pr:classify}.] (a) $\nabla f = (2x + y - 3, \, x + 2y)^\top = \mathbf{0}$ at $(2, -1)$. The Hessian $\begin{pmatrix} 2 & 1 \\ 1 & 2 \end{pmatrix}$ has the eigenvalues $1$ and $3$: a strict local minimum with $f = -3$, which is global (a quadratic function with a positive definite Hessian). (b) $\nabla f = (3x^2 - 3, \, 2y - 2)^\top$, so the critical points are $(\pm 1, 1)$, and $\nabla^2 f = \diag(6x, 2)$. At $(1, 1)$: eigenvalues $6$ and $2$, a strict local minimum with $f = -3$, not global since $f(x, 1) \to -\infty$ as $x \to -\infty$. At $(-1, 1)$: eigenvalues $-6$ and $2$, a saddle point with $f = 1$.
    \item[\ref{pr:four-critical}.] (a) $\nabla f = \big(y(2x + y - 3), \, x(x + 2y - 3)\big)^\top$ vanishes at $(0, 0)$, $(3, 0)$, $(0, 3)$ and $(1, 1)$. With $\nabla^2 f = \begin{pmatrix} 2y & 2x + 2y - 3 \\ 2x + 2y - 3 & 2x \end{pmatrix}$: at the first three points, $\det = -9 < 0$ (saddle points); at $(1, 1)$, the eigenvalues are $1$ and $3$ (a strict local minimum with $f = -1$). (b) The lines $x = 0$, $y = 0$ and $x + y = 3$ bound a triangle on which $f < 0$, and $f$ changes sign across each line. Where two of the lines cross, $f$ has alternating signs in the four sectors, as $xy$ does at the origin, which gives a saddle point. In the triangle, $f$ is negative and $0$ on the boundary, so it has a minimum inside.
    \item[\ref{pr:quadratic-geometry}.] (a) $\nabla^2 f = \begin{pmatrix} 6 & 2 \\ 2 & 6 \end{pmatrix}$: eigenvalue $8$ for $(1, 1)$ and $4$ for $(1, -1)$; indeed $f(t, t) = 8t^2$ and $f(t, -t) = 4t^2$. The level curves are ellipses, longest in the direction $(1, -1)$ of the smaller curvature, with axis ratio $\sqrt{8/4} = \sqrt{2}$. (b) $\nabla^2 g = \begin{pmatrix} 2 & 6 \\ 6 & 2 \end{pmatrix}$: eigenvalue $8$ for $(1, 1)$ and $-4$ for $(1, -1)$, so the origin is a saddle point; $g(t, t) = 8t^2$ increases and $g(t, -t) = -4t^2$ decreases. $g = 0$ on the lines $y = (-3 \pm 2\sqrt{2})x$, i.e., $y \approx -0.17x$ and $y \approx -5.83x$; the other level curves are hyperbolas with these lines as asymptotes.
    \item[\ref{pr:ls-calculus}.] (a) $\nabla L = \frac{1}{3}(28w + 12b - 46, \, 12w + 6b - 20)^\top$; $\nabla L = \mathbf{0}$ gives $14w + 6b = 23$ and $6w + 3b = 10$, so $(w, b) = (\frac{3}{2}, \frac{1}{3})$. (b) $\nabla^2 L = \frac{2}{3}\begin{pmatrix} 14 & 6 \\ 6 & 3 \end{pmatrix}$, with eigenvalues $\frac{17 \pm \sqrt{265}}{3} \approx 11.09$ and $0.240$. Both are positive, and $L$ is a quadratic function with this constant Hessian, so the critical point is the unique global minimum (the example on quadratic functions in the lecture). (c) The residuals are $-\frac{1}{6}, \frac{1}{3}, -\frac{1}{6}$, so $L^* = \frac{1}{18} \approx 0.056$. The normal equations $\Phi^\top\Phi\boldsymbol{\theta} = \Phi^\top\mathbf{y}$ with the rows $(x_i, 1)$ of $\Phi$ are the same two equations.
    \item[\ref{pr:peano}.] (a) $f = y^2 - 3x^2 y + 2x^4$, $\nabla f = (-6xy + 8x^3, \, 2y - 3x^2)^\top$ vanishes at the origin, and $\nabla^2 f(0, 0) = \diag(0, 2)$ has the eigenvalue $0$. (b) On $y = mx$: $f = x^2(m - x)(m - 2x) \approx m^2 x^2$ for $m \neq 0$; on the $x$-axis $f = 2x^4$, and on the $y$-axis $f = y^2$. (c) $f(x, \frac{3}{2}x^2) = \frac{x^2}{2} \cdot (-\frac{x^2}{2}) = -\frac{x^4}{4} < 0$ for $x \neq 0$: between the two parabolas $y = x^2$ and $y = 2x^2$, $f$ is negative.
\end{description}

\section*{Gradient Descent and Newton's Method}

\begin{description}[leftmargin=0pt, itemsep=0.3em]
    \item[\ref{pr:gd-mse}.] (a) $L'(w) = \frac{2}{3}\sum_i (wx_i - y_i)x_i = \frac{2}{3}(14w - 11)$, so $w_{k+1} = w_k - \frac{2\eta}{3}(14w_k - 11)$, and $w^* = \frac{11}{14} \approx 0.786$. (b) $w_1 = 0.367$, $w_2 = 0.562$, $w_3 = 0.667$; the error is multiplied by $1 - \eta \cdot \frac{28}{3} = \frac{8}{15} \approx 0.53$. (c) $|1 - \frac{28}{3}\eta| < 1$, i.e., $0 < \eta < \frac{3}{14} \approx 0.21$; $\eta = \frac{3}{28}$ gives $w_1 = w^*$. (d) $\mathbf{w}_{k+1} = \mathbf{w}_k - \frac{2\eta}{N}X^\top(X\mathbf{w}_k - \mathbf{y}) = \mathbf{w}_k - \frac{2\eta}{N}\sum_i (\mathbf{w}_k^\top\mathbf{x}_i - y_i)\mathbf{x}_i$: each point moves $\mathbf{w}$ along its input $\mathbf{x}_i$, in proportion to its residual (towards fitting it better).
    \item[\ref{pr:gd-quadratic}.] (a) $\nabla f = (4x, 16y)^\top$: $(1, 1) \to (0.6, -0.6) \to (0.36, 0.36)$. (b) The Hessian has the eigenvalues $4$ and $16$, so gradient descent converges if and only if $0 < \eta < \frac{1}{8}$; the best step size is $\eta = \frac{2}{4 + 16} = 0.1$, with $\rho = \frac{\kappa - 1}{\kappa + 1} = 0.6$ for $\kappa = 4$. (c) $0.6^k \leq 10^{-6}$ for $k \geq 28$. (d) $y$ is multiplied by $1 - 16 \cdot 0.125 = -1$ in every step, so it jumps between $\pm y_0$ forever, while $x$ converges.
    \item[\ref{pr:gd-least-squares}.] (a) With the eigenvalues $\lambda_1 \approx 11.09$ and $\lambda_2 \approx 0.240$: convergence for $0 < \eta < \frac{2}{\lambda_1} \approx 0.180$; the best step size is $\frac{2}{\lambda_1 + \lambda_2} = \frac{3}{17} \approx 0.176$, with $\kappa \approx 46.1$ and $\rho \approx 0.958$, so about $319$ steps. (b) The new Hessian is $\frac{2}{3}\begin{pmatrix} \sum(x_i - 2)^2 & \sum(x_i - 2) \\ \sum(x_i - 2) & 3 \end{pmatrix} = \diag(\frac{4}{3}, 2)$, with $\kappa = 1.5$; the best step size is $0.6$, with $\rho = 0.2$, so $9$ steps. Since $w(x - 2) + c = wx + (c - 2w)$, we get $b = c - 2w$ (indeed $c^* = \frac{10}{3}$, $w^* = \frac{3}{2}$, $b^* = \frac{1}{3}$). (c) After centering, the two columns $(x_i - 2)$ and $(1, 1, 1)$ of the design matrix are orthogonal, so the Hessian is diagonal: the slope and the intercept no longer interact, and the level sets are almost circles.
    \item[\ref{pr:newton-log}.] (a) $f'(x) = 1 - \frac{1}{x}$ vanishes at $x = 1$, and $f''(x) = \frac{1}{x^2} > 0$, so $f$ is convex and $f(1) = 1$ is the global minimum. (b) $x - \frac{f'(x)}{f''(x)} = x - (1 - \frac{1}{x})x^2 = 2x - x^2$, and $1 - (2x - x^2) = (1 - x)^2$. From $0.5$: $0.75$, $0.9375$, $0.99609375$, $0.99998474$, with errors $0.25$, $0.0625$, $0.0039$, $1.5 \cdot 10^{-5}$. (c) $1 - x_k = (1 - x_0)^{2^k}$, which tends to $0$ if and only if $|1 - x_0| < 1$, i.e., $0 < x_0 < 2$. For $x_0 = 3$, $x_1 = -3$ is not even in the domain.
    \item[\ref{pr:babylonian}.] $f'(x) = x^2 - 2$ and $f''(x) = 2x$, so $x - \frac{x^2 - 2}{2x} = \frac{1}{2}(x + \frac{2}{x})$. The iterates $1.5$, $1.416667$, $1.414216$ converge to $\sqrt{2}$, with errors $0.086$, $0.0025$, $2.1 \cdot 10^{-6}$. For $\sqrt{a}$, minimize $\frac{x^3}{3} - ax$, which gives $x_{k+1} = \frac{1}{2}(x_k + \frac{a}{x_k})$.
    \item[\ref{pr:smooth-descent}.] (a) $s' = \sigma$ and $s'' = \sigma(1 - \sigma) \in (0, \frac{1}{4}]$. (b) With $\eta = \frac{1}{L} = 4$: $s(x_{k+1}) \leq s(x_k) - 2\sigma(x_k)^2$, and $\min_{k < K}\sigma(x_k)^2 \leq \frac{s(x_0)}{2K}$, since $s > 0$. The iterates $x_{k+1} = x_k - 4\sigma(x_k)$ decrease to $-\infty$, $s(x_k) \to 0 = \inf s$, and the gradients $\sigma(x_k) \to 0$. The theorem only promises small gradients, not a minimum, so this is consistent.
\end{description}

\section*{Convexity}

\begin{description}[leftmargin=0pt, itemsep=0.3em]
    \item[\ref{pr:convex-which}.] (a) Convex, even strictly: the Hessian has the eigenvalues $1$ and $3$. (b) Not convex: $\begin{pmatrix} 2 & 4 \\ 4 & 2 \end{pmatrix}$ has the eigenvalue $-2$. (c) Convex: the Hessian $e^{\norm{\mathbf{x}}^2}(2I + 4\mathbf{x}\mathbf{x}^\top)$ is positive definite. (d) Convex, like every norm. (e) Convex: $f''(x) = \frac{1}{4}x^{-3/2} > 0$. (f) Convex: the Hessian is $\frac{2}{y^3}\begin{pmatrix} y^2 & -xy \\ -xy & x^2 \end{pmatrix} = \frac{2}{y^3}\mathbf{v}\mathbf{v}^\top$ with $\mathbf{v} = (y, -x)^\top$, which is positive semidefinite. (g) Not convex: the Hessian $\begin{pmatrix} 0 & -1/y^2 \\ -1/y^2 & 2x/y^3 \end{pmatrix}$ has determinant $-\frac{1}{y^4} < 0$.
    \item[\ref{pr:am-gm}.] (a) $(\log x)'' = -\frac{1}{x^2} < 0$. (b) Concavity gives $\log(\theta a + (1 - \theta)b) \geq \theta\log a + (1 - \theta)\log b$; apply $\exp$, which is increasing. (c) With $a = x$, $b = \frac{4}{x}$ and $\theta = \frac{1}{2}$: $x + \frac{4}{x} \geq 2\sqrt{4} = 4$, with equality for $x = \frac{4}{x}$, i.e., $x = 2$. Alternatively, the derivative $1 - \frac{4}{x^2}$ vanishes at $x = 2$, and the second derivative $\frac{8}{x^3}$ is positive, so the function is convex on $(0, \infty)$ and $2$ is the global minimum.
    \item[\ref{pr:convex-losses}.] (a) Differentiating $\nabla\ell = -y\,\sigma(-y\,\mathbf{w}^\top\mathbf{x})\,\mathbf{x}$ gives $\nabla^2\ell = \sigma(1 - \sigma)\,y^2\,\mathbf{x}\mathbf{x}^\top = \sigma(1 - \sigma)\,\mathbf{x}\mathbf{x}^\top$ (with $\sigma$ evaluated at $-y\,\mathbf{w}^\top\mathbf{x}$), which is positive semidefinite. (b) It is the maximum of the affine functions $0$ and $1 - y\,\mathbf{w}^\top\mathbf{x}$. (c) $\norm{X\mathbf{w} - \mathbf{y}}^2$ has the positive semidefinite Hessian $2X^\top X$, $\lambda\norm{\mathbf{w}}_1$ is a nonnegative multiple of a norm, and sums of convex functions are convex. (d) $f(0) = 0$ and $f(2) = 1$, but $f(1) = 1 > \frac{1}{2}f(0) + \frac{1}{2}f(2)$.
    \item[\ref{pr:first-order-optimality}.] $f(\mathbf{y}) \geq f(\mathbf{x}^*) + \nabla f(\mathbf{x}^*)^\top(\mathbf{y} - \mathbf{x}^*) \geq f(\mathbf{x}^*)$ for all $\mathbf{y} \in C$, by the first-order characterization of convexity. Picture: at a boundary point $\mathbf{x}^*$ of a half-plane, the steepest descent direction $-\nabla f(\mathbf{x}^*)$ points straight out of $C$, so we cannot decrease $f$ without leaving $C$. (This is the KKT condition for one linear constraint.)
    \item[\ref{pr:two-layer-scalar}.] (a) $f \geq 0$, and $f = 0$ exactly on the hyperbola $w_1 w_2 = 1$. (b) $(1, 1)$ and $(-1, -1)$ are minima with $f = 0$, but their midpoint $(0, 0)$ has $f = 1$. (c) $\nabla f = 2(w_1 w_2 - 1)(w_2, w_1)^\top$ vanishes on the hyperbola and at the origin. At the origin, $\nabla^2 f = \begin{pmatrix} 0 & -2 \\ -2 & 0 \end{pmatrix}$ has the eigenvalues $\pm 2$: a saddle point. At $(1, 1)$, $\nabla^2 f = \begin{pmatrix} 2 & 2 \\ 2 & 2 \end{pmatrix}$ has the eigenvalues $4$ (for $(1, 1)$) and $0$ (for $(1, -1)$, the tangent direction of the hyperbola, along which $f$ stays $0$).
    \item[\ref{pr:strong-convexity}.] (a) The Hessian $\begin{pmatrix} 2 & 1 \\ 1 & 2 \end{pmatrix}$ has the eigenvalues $1$ and $3$: $\mu = 1$, $L = 3$. (b) $(1 - \frac{1}{3})^k \leq 10^{-6}$ for $k \geq 35$. (c) $f(1, 0) = 1$, $f^* = 0$, $\nabla f(1, 0) = (2, 1)^\top$, and $\frac{1}{2}\norm{\nabla f}^2 = 2.5 \geq 1$.
    \item[\ref{pr:subgradient-practice}.] $f'(x) = -3$ for $x < -1$, $1$ for $-1 < x < 1$, and $3$ for $x > 1$; $\partial f(-1) = [-3, 1]$ and $\partial f(1) = [1, 3]$. Only $\partial f(-1)$ contains $0$, so the minimum is at $x^* = -1$, with $f(-1) = 2$.
    \item[\ref{pr:soft-threshold-practice}.] (a) $(1.5, 0, -0.7, 0)^\top$. (b) $w^* = S_2(3) = 1$; the condition is $0 \in (1 - 3) + 2\,\partial|\cdot|(1) = \{0\}$. The minimum is at $0$ if and only if $|a| \leq 2$.
    \item[\ref{pr:group-lasso}.] (a) For $\mathbf{x} \neq \mathbf{0}$, the norm is differentiable with the gradient $\mathbf{x}/\norm{\mathbf{x}}$. At $\mathbf{0}$, $\norm{\mathbf{y}} \geq \mathbf{g}^\top\mathbf{y}$ for all $\mathbf{y}$ holds if $\norm{\mathbf{g}} \leq 1$ (Cauchy--Schwarz) and fails for $\mathbf{y} = \mathbf{g}$ if $\norm{\mathbf{g}} > 1$. (b) Check the criterion $\frac{1}{\tau}(\mathbf{v} - \mathbf{p}) \in \partial\norm{\cdot}(\mathbf{p})$: for $\norm{\mathbf{v}} > \tau$, $\mathbf{p} = (1 - \frac{\tau}{\norm{\mathbf{v}}})\mathbf{v}$ and $\frac{1}{\tau}(\mathbf{v} - \mathbf{p}) = \frac{\mathbf{v}}{\norm{\mathbf{v}}} = \frac{\mathbf{p}}{\norm{\mathbf{p}}}$; for $\norm{\mathbf{v}} \leq \tau$, $\mathbf{p} = \mathbf{0}$ and $\norm{\frac{1}{\tau}\mathbf{v}} \leq 1$.
\end{description}

\section*{Constrained Optimization}

\begin{description}[leftmargin=0pt, itemsep=0.3em]
    \item[\ref{pr:lagrange-ellipse}.] $(y, x) + \lambda(2x, 8y) = (0, 0)$ gives $y = -2\lambda x$ and $x = -8\lambda y = 16\lambda^2 x$, so $\lambda = \pm\frac{1}{4}$ (the case $x = 0$ forces $y = 0$, which is not on the ellipse). For $\lambda = -\frac{1}{4}$: $y = \frac{x}{2}$, $x = \pm 2$, so the maximum is $2$ at $\pm(2, 1)$. For $\lambda = \frac{1}{4}$: the minimum is $-2$ at $\pm(2, -1)$. At $(2, 1)$, $\nabla(xy) = (1, 2)$ and $\nabla g = (4, 8)$ are parallel, and both curves have the tangent line $x + 2y = 4$.
    \item[\ref{pr:max-xyz}.] (a) $(yz, xz, xy) = \mu(1, 1, 1)$ gives $x = y = z = 4$, and the maximum is $64$. (b) The Lagrange conditions $y + 2z = \mu yz$, $x + 2z = \mu xz$, $2x + 2y = \mu xy$ give $x = y$ and then $x = 2z$, so $4z^3 = 32$: $x = y = 4$, $z = 2$, with surface area $48$.
    \item[\ref{pr:distance-ellipse}.] Optimize $\norm{\mathbf{x}}^2$ subject to $\mathbf{x}^\top Q\mathbf{x} = 3$ with $Q = \begin{pmatrix} 1 & 1/2 \\ 1/2 & 1 \end{pmatrix}$. The Lagrange condition $2\mathbf{x} - 2\mu Q\mathbf{x} = \mathbf{0}$ says that $\mathbf{x}$ is an eigenvector of $Q$, and then $\norm{\mathbf{x}}^2 = \frac{3}{\text{eigenvalue}}$. Eigenvalue $\frac{3}{2}$ for $(1, 1)$: the closest points $\pm(1, 1)$, at distance $\sqrt{2}$. Eigenvalue $\frac{1}{2}$ for $(1, -1)$: the farthest points $\pm(\sqrt{3}, -\sqrt{3})$, at distance $\sqrt{6}$.
    \item[\ref{pr:min-norm-practice}.] $A = \begin{pmatrix} 1 & 1 & 1 \\ 1 & 0 & -1 \end{pmatrix}$, $AA^\top = \diag(3, 2)$, $(AA^\top)^{-1}\mathbf{b} = (1, \frac{1}{2})^\top$, and $\mathbf{x}^* = A^\top(1, \frac{1}{2})^\top = (\frac{3}{2}, 1, \frac{1}{2})^\top$, with $\norm{\mathbf{x}^*}^2 = \frac{7}{2}$. The null space is spanned by $\mathbf{v} = (1, -2, 1)^\top$, $\mathbf{x}^{*\top}\mathbf{v} = 0$, and $\norm{\mathbf{x}^* + t\mathbf{v}}^2 = \frac{7}{2} + 6t^2$.
    \item[\ref{pr:kkt-projection}.] (a) Stationarity $2(x - 3) + \mu = 0$, $2(y - 2) + \mu = 0$. For $\mu = 0$, the point $(3, 2)$ is not feasible, so $x + y = 2$, which gives $\mu = 3 \geq 0$ and $(x, y) = (1.5, 0.5)$; this is the global minimum, since the problem is convex. It is the projection of $(3, 2)$ onto the half-plane, at distance $\frac{3}{\sqrt{2}} = \frac{|3 + 2 - 2|}{\sqrt{2}}$. (b) The minimum is $(0, 2)$. The constraints $x + y \leq 2$ and $-x \leq 0$ are active, with $\mu_1 = \mu_2 = 6$, and $-y \leq 0$ is inactive, with $\mu_3 = 0$: indeed $\nabla f(0, 2) = (0, -6)$ and $(0, -6) + 6(1, 1) + 6(-1, 0) = (0, 0)$.
    \item[\ref{pr:box-sphere}.] With the half-sides $x, y, z > 0$, maximize $8xyz$ subject to $x^2 + y^2 + z^2 = 1$. The Lagrange conditions give $x = y = z = \frac{1}{\sqrt{3}}$: a cube with side $\frac{2}{\sqrt{3}} \approx 1.155$ and volume $\frac{8}{3\sqrt{3}} \approx 1.540$, which is the fraction $\frac{2}{\sqrt{3}\,\pi} \approx 36.8\%$ of the volume $\frac{4\pi}{3}$ of the ball.
    \item[\ref{pr:dual-practice}.] (a) $g(\mathbf{x}) = 1 - x_1 - x_2$. The Lagrangian $\norm{\mathbf{x}}^2 + \mu(1 - x_1 - x_2)$ is smallest at $\mathbf{x} = \frac{\mu}{2}(1, 1)$, so $d(\mu) = \mu - \frac{\mu^2}{2}$. (b) $\mu^* = 1$ and $d^* = \frac{1}{2}$. KKT: $\mathbf{x}^* = (\frac{1}{2}, \frac{1}{2})$ with $\mu = 1$, and $p^* = \frac{1}{2}$.
    \item[\ref{pr:duality-gap}.] (a) (A): $\mathcal{L} = (\mu - 1)x^2 - \mu$, so $d(\mu) = -\infty$ for $\mu < 1$ and $d(\mu) = -\mu$ for $\mu \geq 1$; $d^* = -1$. (B): $\mathcal{L} = -x^2 + (\mu_1 - \mu_2)x - \mu_1 - \mu_2$ is unbounded below, so $d \equiv -\infty$ and $d^* = -\infty$. (b) Only for (A). The objective $-x^2$ is not convex, so the theorem does not apply; whether strong duality holds then depends on how the constraints are written.
    \item[\ref{pr:svm-dual-practice}.] (a) Maximize $\mu_1 + \mu_2 + \mu_3 - \frac{1}{2}\norm{(\mu_1 + \mu_2 + 2\mu_3)(1, 1)}^2 = \mu_1 + \mu_2 + \mu_3 - (\mu_1 + \mu_2 + 2\mu_3)^2$ subject to $\boldsymbol{\mu} \geq \mathbf{0}$ and $\mu_1 - \mu_2 + \mu_3 = 0$. (b) $\mu_1 - \mu_2 + \mu_3 = 0$; the dual value is $\frac{1}{2} - \frac{1}{4} = \frac{1}{4}$; $\mathbf{w} = (\frac{1}{2}, \frac{1}{2})^\top$, and $b = 0$ from $y_1(\mathbf{w}^\top\mathbf{x}_1 + b) = 1$. (c) $y_i(\mathbf{w}^\top\mathbf{x}_i + b) = 1, 1, 2 \geq 1$, and $\frac{1}{2}\norm{\mathbf{w}}^2 = \frac{1}{4}$. By weak duality, no feasible point has a smaller value, and no $\boldsymbol{\mu}$ a larger one. The margin is $\frac{1}{\norm{\mathbf{w}}} = \sqrt{2}$, and the support vectors are $\mathbf{x}_1$ and $\mathbf{x}_2$.
\end{description}

\section*{Vector Fields: Divergence and Curl}

\begin{description}[leftmargin=0pt, itemsep=0.3em]
    \item[\ref{pr:div-curl}.] (a) $\Div F = 2x + 2x = 4x$, $\curl F = 2y - (-2y) = 4y$; not a gradient field. (b) $\Div F = 2 - 3 = -1$, $\curl F = 1 - 1 = 0$; $f = x^2 + xy - \frac{3}{2}y^2$. (c) $\Div F = -(x^2 + y^2)\sin(xy)$, and $\curl F = (\cos(xy) - xy\sin(xy)) - (\cos(xy) - xy\sin(xy)) = 0$; $f = \sin(xy)$. (d) $\Div F = 0 + 2x + 0 = 2x$, $\curl F = (1 - 1, \, 0 - 0, \, 2y - 2y) = \mathbf{0}$; $f = xy^2 + yz$.
    \item[\ref{pr:laplacian}.] (a) $2 + 2 - 4 = 0$; $e^x\sin y - e^x\sin y = 0$; $2n$. The first two are harmonic. (b) $\Delta f = \tr\nabla^2 f = \lambda_1 + \lambda_2 = 0$, so the eigenvalues of the Hessian cannot both be positive. (c) $\mathbf{s}(\mathbf{x}) = -\Sigma^{-1}(\mathbf{x} - \boldsymbol{\mu})$, with the constant Jacobian $-\Sigma^{-1}$, so $\Delta\log p = -\tr\Sigma^{-1}$, which is $-(1 + \frac{1}{4}) = -1.25$ for $\Sigma = \diag(1, 4)$. The score field is a sink everywhere: it points towards the mean.
    \item[\ref{pr:stretch-rotate}.] $S = \left(\begin{smallmatrix} 1 & 1 \\ 1 & 1 \end{smallmatrix}\right)$, with the eigenvalue $2$ for $(1, 1)$ and $0$ for $(1, -1)$, and $W = \left(\begin{smallmatrix} 0 & 2 \\ -2 & 0 \end{smallmatrix}\right)$, so $\curl F = -1 - 3 = -4$, and $\Div F = \tr A = 2$, the sum of the stretching rates. The field stretches along $(1, 1)$ at the rate $2$, does not stretch along $(1, -1)$, and rotates clockwise with the angular velocity $\frac{1}{2}|\curl F| = 2$. (At $(1, 0)$, the arrow is $(1, -1)$, which points clockwise around the origin. In the section on ordinary differential equations, we will see that the solutions of $\mathbf{x}' = A\mathbf{x}$ spiral outwards, since the eigenvalues $1 \pm i\sqrt{3}$ of $A$ have a positive real part.)
    \item[\ref{pr:gda}.] (a) Descent in $x$ follows $-\frac{\partial}{\partial x}(xy) = -y$, and ascent in $y$ follows $\frac{\partial}{\partial y}(xy) = x$. $\Div F = 0$ and $\curl F = 2$. A field of the form $-\nabla g$ has zero curl, since the Hessian of $g$ is symmetric, so $F$ is not of this form. (b) $(x - \eta y)^2 + (y + \eta x)^2 = (1 + \eta^2)(x^2 + y^2)$. (c) With $x_+ = x - \eta y$ and $y_+ = y + \eta x_+ = \eta x + (1 - \eta^2)y$, expanding gives $x_+^2 - \eta x_+y_+ + y_+^2 = x^2 - \eta xy + y^2$. For $0 < \eta < 2$, the matrix $\left(\begin{smallmatrix} 1 & -\eta/2 \\ -\eta/2 & 1 \end{smallmatrix}\right)$ of $Q$ has the positive eigenvalues $1 \pm \frac{\eta}{2}$, so the level sets of $Q$ are ellipses, and the iterates stay on one of them.
\end{description}

\section*{Multiple Integrals and the Divergence Theorem}

\begin{description}[leftmargin=0pt, itemsep=0.3em]
    \item[\ref{pr:fubini-practice}.] (a) The region is the triangle $0 \leq y \leq x \leq 1$, so the integral is $\int_0^1\int_0^x e^{x^2}\, dy\, dx = \int_0^1 xe^{x^2}\, dx = \frac{e - 1}{2} \approx 0.859$. (b) $\int_0^{\pi/2}\int_0^1 r^2\cos\varphi\sin\varphi \cdot r\, dr\, d\varphi = \frac{1}{4} \cdot \frac{1}{2} = \frac{1}{8}$. (c) $\frac{1}{3} + 2 \cdot \frac{1}{4} + \frac{1}{3} = \frac{7}{6}$.
    \item[\ref{pr:prob-integrals}.] (a) The area of the triangle $x, y \geq 0$, $x + y \leq \frac{1}{2}$, which is $\frac{1}{8}$. (b) For $x \leq \frac{1}{2}$, every $y$ works; for $x > \frac{1}{2}$, $y \leq \frac{1}{2x}$. So the probability is $\frac{1}{2} + \int_{1/2}^1 \frac{dx}{2x} = \frac{1}{2} + \frac{\log 2}{2} \approx 0.847$. (c) By symmetry, $2\int_0^1\int_0^x x\, dy\, dx = 2\int_0^1 x^2\, dx = \frac{2}{3}$.
    \item[\ref{pr:gauss-ellipse}.] (a) $\mathbf{x}^\top A\mathbf{x} = \norm{B\mathbf{x}}^2$, and $\mathbf{x} = B^{-1}\mathbf{u}$ has the Jacobian determinant $\det B^{-1}$, with $|\det B| = \sqrt{\det A}$. So the area is $\frac{1}{\sqrt{\det A}}$ times the area $\pi$ of the unit disk, and the integral is $\frac{1}{\sqrt{\det A}}$ times $\int e^{-\norm{\mathbf{u}}^2/2}\, d\mathbf{u} = 2\pi$. (b) $A = \left(\begin{smallmatrix} 2 & 1 \\ 1 & 1 \end{smallmatrix}\right)$ has $\det A = 1$: the area is $\pi$ and the integral is $2\pi$.
    \item[\ref{pr:line-integrals}.] (a) $F$: along $\gamma_1$, $\int_0^1 (-t^2 \cdot 1 + t \cdot 2t)\, dt = \frac{1}{3}$; along $\gamma_2$, $\int_0^1 (-t + t)\, dt = 0$. $G$: $\int_0^1 (t^2 + 2t^2)\, dt = 1$ and $\int_0^1 2t\, dt = 1$. $G$ has zero curl and is a gradient field on $\mathbb{R}^2$, so its line integrals only depend on the endpoints; $F$ has $\curl F = 2$. (b) $f = xy$, and $f(1, 1) - f(0, 0) = 1$. (c) The closed curve goes counterclockwise around the region between $y = x^2$ and $y = x$, of area $\int_0^1 (t - t^2)\, dt = \frac{1}{6}$. The circulation is $\frac{1}{3} - 0 = \frac{1}{3} = 2 \cdot \frac{1}{6}$, the integral of the curl.
    \item[\ref{pr:score-matching-1d}.] (a) $\frac{1}{2}s_\theta^2 = \frac{1}{2}\theta^2x^2$ and $\Div s_\theta = s_\theta' = -\theta$. (b) Minimize $\frac{1}{2}\theta^2 m_2 - \theta$ with $m_2 = \frac{1}{N}\sum_i x_i^2$: $\hat\theta = \frac{1}{m_2}$, the same as the MLE. (c) $m_2 = \frac{1 + 4 + 9 + 0 + 4}{5} = 3.6$, so $\hat\theta = \frac{1}{3.6} \approx 0.278$.
    \item[\ref{pr:score-function}.] (a) With $g = x^2$, the formula gives $g'(\theta) = \int x^2(x - \theta)p_\theta(x)\, dx = \int (\theta + z)^2z\,\varphi(z)\, dz = \int (\theta^2z + 2\theta z^2 + z^3)\varphi(z)\, dz = 2\theta$, since the odd moments vanish and $\int z^2\varphi = 1$. (b) $g'(\theta) = \int 2(\theta + z)\varphi(z)\, dz = 2\theta$. (c) Both estimates have mean $0$; the mean squares are $\int z^6\varphi = 15$ for $X^3$ and $4\int z^2\varphi = 4$ for $2X$. The reparameterization estimator is less noisy, which is why it is preferred when it is available (e.g., in variational autoencoders).
    \item[\ref{pr:coulomb}.] (a) $\frac{\partial}{\partial x_i}\frac{x_i}{\norm{\mathbf{x}}^3} = \frac{1}{\norm{\mathbf{x}}^3} - \frac{3x_i^2}{\norm{\mathbf{x}}^5}$, and the sum over $i$ is $\frac{3}{\norm{\mathbf{x}}^3} - \frac{3}{\norm{\mathbf{x}}^3} = 0$. (b) On the sphere, $\mathbf{n} = \frac{\mathbf{x}}{R}$ and $F^\top\mathbf{n} = \frac{1}{R^2}$, so the flux is $\frac{4\pi R^2}{R^2} = 4\pi$. (c) $F$ is not defined at the origin, so the theorem does not apply to the ball. For the region between two spheres, it applies: the flux out of the outer sphere is $4\pi$, the flux through the inner sphere (whose outward normal points towards the origin) is $-4\pi$, and the total is $0 = \int \Div F$. The point charge is a source of total strength $4\pi$ concentrated at the origin.
\end{description}

\section*{Ordinary Differential Equations}

\begin{description}[leftmargin=0pt, itemsep=0.3em]
    \item[\ref{pr:scalar-ode}.] (a) $x(t) = 3e^{-2t}$, $x(1) = 3e^{-2} \approx 0.406$; it halves after $\frac{\log 2}{2} \approx 0.347$. (b) With $u = 4e^{-t}$: $x' = \frac{u}{(1 + u)^2}$ and $x(1 - x) = \frac{1}{1 + u}\cdot\frac{u}{1 + u}$, which agree, and $x(0) = \frac{1}{5}$. So $x(t) = \sigma(t - \log 4)$, $x(2) \approx 0.649$, and $x(t) \to 1$.
    \item[\ref{pr:existence}.] (a) $x' = \frac{1}{(1 - t)^2} = x^2$ and $x(0) = 1$. The solution tends to $\infty$ as $t \to 1$: it blows up in finite time. The theorem does not apply, since $F(x) = x^2$ is not Lipschitz continuous on $\mathbb{R}$: $|x^2 - y^2| = |x + y||x - y|$, and $|x + y|$ is unbounded. (b) For $x = t^2$: $x' = 2t = 2\sqrt{t^2}$. The function $2\sqrt{|x|}$ is not Lipschitz continuous near $0$, where its slope is unbounded, so uniqueness can fail.
    \item[\ref{pr:linear-system}.] (a) Eigenvalue $1$ with eigenvector $(1, 1)$, and $-3$ with $(1, -1)$; $(1, 0) = \frac{1}{2}(1, 1) + \frac{1}{2}(1, -1)$, so $\mathbf{x}(t) = \frac{1}{2}e^t(1, 1)^\top + \frac{1}{2}e^{-3t}(1, -1)^\top$ and $\mathbf{x}(1) \approx (1.384, 1.334)^\top$. (b) Exactly the multiples of $(1, -1)$, i.e., the line $x_2 = -x_1$; all other solutions grow along $(1, 1)$. (c) $e^{At} = \frac{1}{2}\begin{pmatrix} e^t + e^{-3t} & e^t - e^{-3t} \\ e^t - e^{-3t} & e^t + e^{-3t} \end{pmatrix}$.
    \item[\ref{pr:rotation-ode}.] (a) $\frac{d}{dt}\norm{\mathbf{x}}^2 = 2\mathbf{x}^\top A\mathbf{x} = 0$, and $e^{At}$ is the rotation by the angle $\omega t$, so the period is $\frac{2\pi}{\omega}$. (b) $\norm{\mathbf{x} + hA\mathbf{x}}^2 = \norm{\mathbf{x}}^2 + 2h\mathbf{x}^\top A\mathbf{x} + h^2\norm{A\mathbf{x}}^2 = (1 + h^2\omega^2)\norm{\mathbf{x}}^2$. One period takes $\frac{2\pi}{0.1} \approx 63$ steps, so the norm grows by $(1.01)^{31.4} \approx 1.37$.
    \item[\ref{pr:gradient-flow}.] (a) $x_1' = -2x_1$ and $x_2' = -8x_2$, so $\mathbf{x}(t) = (2e^{-2t}, e^{-8t})$ and $f(\mathbf{x}(t)) = 4e^{-4t} + 4e^{-16t}$, which decreases. (b) $\mathbf{x}_{k+1} = \mathbf{x}_k - h\nabla f(\mathbf{x}_k) = \big((1 - 2h)x_{1,k}, (1 - 8h)x_{2,k}\big)$, which is gradient descent with $\eta = h$. It converges if $|1 - 2h| < 1$ and $|1 - 8h| < 1$, i.e., $0 < h < 0.25$. For $h = 0.2$, the factors are $0.6$ and $-0.6$ (convergence with oscillation in $x_2$); for $h = 0.3$, $x_2$ is multiplied by $-1.4$ and diverges, although the exact flow converges.
    \item[\ref{pr:euler-error}.] $0.75^4 \approx 0.316$ and $0.875^8 \approx 0.344$, compared with $e^{-1} \approx 0.368$; the errors $0.051$ and $0.024$ differ by a factor of about $2$, as expected for a method with error proportional to $h$.
\end{description}

\section*{Flows, Volumes and Densities}

\begin{description}[leftmargin=0pt, itemsep=0.3em]
    \item[\ref{pr:volume-flow}.] (a) $\Div F = \tr A = -2$, so areas are multiplied by $e^{-2t}$. The eigenvalues are $1$ and $-3$ (triangular matrix), and $\det e^{At} = e^te^{-3t} = e^{-2t}$. (b) $DF = \diag(1 - 3x^2, -1)$, so $\Div F = -3x^2$. $F = \mathbf{0}$ at $(0, 0)$ (divergence $0$) and $(\pm 1, 0)$ (divergence $-3$). Areas are preserved to first order on the line $x = 0$ and compressed everywhere else.
    \item[\ref{pr:continuity-1d}.] (a) $P(x(t) \leq a) = P(x(0) \leq e^ta)$; differentiating in $a$ gives $e^t\varphi(e^ta)$. (b) With $\varphi'(u) = -u\varphi(u)$: $\frac{\partial p_t}{\partial t} = e^t\varphi(e^tx) + e^{2t}x\,\varphi'(e^tx)$, and $-\frac{\partial}{\partial x}(-xp_t) = p_t + x\frac{\partial p_t}{\partial x} = e^t\varphi(e^tx) + xe^{2t}\varphi'(e^tx)$. (c) $\log p_t(x(t)) = t + \log\varphi(x(0))$, which increases at the rate $1$. The flow contracts lengths by $e^{-t}$ (here $\Div F = -1$), so the same probability is squeezed into a smaller interval, and the density increases.
\end{description}

\end{document}
